% 本文件由 LaTeX 排版 Agent 根据有效 translation.json 自动生成。
% author=Lactantius; work_id=0706; title=拉克坦提乌斯残篇

\chapter{拉克坦提乌斯残篇}

% structure: p0001 source=h1 target=omitted reason=page-title-already-rendered-by-parent

一、恐惧、爱、喜乐、悲伤、欲念、热望、愤怒、怜悯、好胜、惊佩——这些心灵的悸动或情感从主创造人之初就存在；它们被有益而有利地引入人性，使人能够通过有条理地并按理性支配自己，以勇敢的行动操练那些善德，并借此堪当从主领受永生。因为这些心灵的情感若被约束在适当界限内，即正确运用，便会结出今生的善德和未来的永赏。但若它们越出界限，即转向邪恶之路，那么恶习和不义便产生，并带来永罚。\par

二、据我们所忆，\Person{拉克坦提乌斯}也论及诗律——即\OriginalTerm{五音步诗行}{pentameter}和\OriginalTerm{四音步诗行}{tetrameter}，如他所说。\par

三、\Person{菲尔米亚努斯}写信给\Person{普罗布斯}讨论喜剧的诗律，如此说道：\Quote{\Person{泰伦提乌斯}至于你提出的关于喜剧诗律的问题，我也知道许多人尤其认为\Person{泰伦提乌斯}的剧本不符合希腊喜剧的诗律——即\Person{米南德}、\Person{菲勒蒙}和\Person{狄菲洛斯}的剧作，它们由\OriginalTerm{三音步}{trimeter}诗句组成；因为我们的古代喜剧作家在剧本的韵律调制上更喜欢追随\Person{欧波利斯}、\Person{克拉提努斯}和\Person{阿里斯托芬}，如前所述。} 在\Person{泰伦提乌斯}和\Person{普劳乌斯}的剧本中，以及其他喜剧和悲剧作家的作品中，有一种度量——即诗律——存在，让以下这些人来证明：\Person{西塞罗}、\Person{斯考鲁斯}和\Person{菲尔米亚努斯}。\par

四、我们将引用我们的\Person{拉克坦提乌斯}就此问题在其致\Person{普罗布斯}的第三卷中所表达的见解。他说，高卢人自古以来就被称为加拉达人，因其身体白皙；\Person{西彼拉}女先知亦如此称呼他们。而诗人说以下诗句时正是此意——\par

\Quote{金项圈装饰着他们乳白色的颈项，}\par

而他本可使用\Emph{白色}一词。显然，由此行省得名加拉达，高卢人抵达该地后与希腊人结合，因而该地被称为高卢希腊，后称加拉达。他如此论及加拉达人，并记述西方一民族跨越如此遥远的中土距离，定居于东方一地区，这并不奇怪。\par

\SourceRecord{Translated by William Fletcher.；Ante-Nicene Fathers；Vol. 7.；Edited by Alexander Roberts, James Donaldson, and A. Cleveland Coxe.；Buffalo, NY: Christian Literature Publishing Co.,；1886.；<http://www.newadvent.org/fathers/0706.htm>.}
